% ============================================================
% 第 1 章 问题重述
% ============================================================

\section{问题重述}

\subsection{问题背景}

啤酒的风味、口感与香气由其化学组成决定。发酵过程中，酿酒酵母的糖代谢与酯类、高级醇合成，酒花添加带来的萜类化合物，以及麦芽和微生物代谢产生的醛酮类物质，共同塑造了啤酒的果香、醇厚感、苦味与纯净度\citep{vanderhaegen2006chemistry}。同一化合物在不同浓度区间内可能呈现完全不同的感官贡献；多种化合物之间又存在显著共表达或共同代谢关系。因此，仅靠单变量分析难以刻画“化合物谱—感官画像”的整体联系。

附件给出了 250 个商用啤酒样本的化合物丰度、化合物相关矩阵、感官评分卡、感官相关矩阵以及综合参考数据。数据同时具有高维、小样本、缺失、异常值与多源异质的特点：化合物丰度部分列经过对数变换、部分列保留线性浓度，感官评分以 0--6 离散等级记录，部分表格只给出标准化后的连续值，跨表连接还存在标识列不一致的问题。上述特点决定了本文必须先完成严格的数据对齐与预处理，再针对三问分别设计模型。

\subsection{数据说明}

附件共 8 份数据，其内容与本文用途如表~\ref{tab:data-overview} 所示。

\begin{table}[htbp]
  \centering
  \caption{附件数据一览与本文用途}
  \label{tab:data-overview}
  \small
  \begin{tabularx}{\textwidth}{@{}l X X@{}}
    \toprule
    文件 & 主要内容 & 本文用途 \\
    \midrule
    \texttt{1.csv} & 250 个啤酒样本 $\times$ 化合物丰度矩阵；含少量理化参数，化合物列逐列核对（37 列 log10、33 列线性、161 列无校准按 log10）。 & 问题一至问题三的主特征表。 \\
    \texttt{2.csv} & 化合物两两之间的 Pearson 相关矩阵。 & 问题二共表达图的边权来源。 \\
    \texttt{3.xlsx} & Aroma 与 Flavour 两套感官评分卡，评分档位为 0--6。 & 问题一的目标维度定义。 \\
    \texttt{4.csv} & 感官评分的标准化/中心化连续值。 & 问题一、二的监督信号。 \\
    \texttt{5.csv} & 感官属性之间的相关矩阵。 & 感官维度结构分析与共线性诊断。 \\
    \texttt{6.csv} & 化合物/理化指标与感官属性的交叉相关矩阵。 & 问题二的外部一致性校验。 \\
    \texttt{7.csv} & 250 个样本的感官描述符，含风格标签 \texttt{tasting\_category\_fine}。 & 问题三的风格标签；问题一的辅助监督。 \\
    \texttt{8.xlsx} & 感官阈值、感官名称、GBM/SHAP、校准化学数据、价格等。 & 阈值参照、变换核验与外部校验。 \\
    \bottomrule
  \end{tabularx}
\end{table}

需要特别说明两点。第一，\texttt{4.csv} 与 \texttt{7.csv} 的 \texttt{beer\_id} 存在科学计数法存储，唯一值数量与样本数不一致；本文跨表连接统一使用 \texttt{beer} 列，不使用 \texttt{beer\_id}。第二，\texttt{3.xlsx} 的 beer 评分卡是空白评分卡而非已填写的 0--6 分数，因此问题一不存在可直接监督的等级标签；本文把 \texttt{4.csv}、\texttt{7.csv} 的标准化连续值作为连续监督信号，再通过有序校准得到 0--6 分级结果，并在结果中明确这一限制。

\subsection{问题提出}

\begin{enumerate}
  \item \textbf{问题一：数据缺失条件下的感官评分。}以 \texttt{1.csv} 的化合物数据为基础，给出 \texttt{3.xlsx} 中定义的 Aroma 与 Flavour 两套、共约 40 个感官维度的 0--6 分制分级评分结果。数据存在缺失值与异常值，且没有现成的等级标签可供直接训练。
  \item \textbf{问题二：化合物—风味分析。}利用 \texttt{2.csv} 的化合物相关数据构建化合物共表达图，结合 \texttt{1.csv} 特征，给出每个感官维度的 Top-3 关键因素化合物，并从 bootstrap 重采样、不同建模方法一致性以及化学代谢路径三个角度论证结果的稳定性与合理性。
  \item \textbf{问题三：风味与“风味指纹”生成。}以 \texttt{tasting\_category\_fine} 为风格标签，学习风格条件生成模型：给定风格（如 Tripel、Kriek、Lambic）生成合理的化合物谱与感官画像；并完成跨风格迁移（如把 Pils/Lager 的化合物谱迁移为 Tripel 风格），同时保证 DMS、双乙酰、4VG、反式-2-壬烯醛、stale hops 等缺陷化合物不恶化。
\end{enumerate}

\subsection{建模路线与逻辑递进}

三问共享同一套数据底座与符号体系，但任务性质不同：问题一是“从化合物到感官”的有监督映射，问题二是“从相关结构到关键因素”的关联分析，问题三是“从风格标签到化合物谱”的条件生成。为避免一上来堆砌复杂模型，本文对每一问都按“基线 $\rightarrow$ 改进 $\rightarrow$ 最终”组织：先给一个可复现、可解释的简单基线，再针对基线的具体缺陷逐层改进，最后用消融与稳定性实验说明每一层改进的必要性。总体路线见表~\ref{tab:roadmap} 与图~\ref{fig:roadmap}。

\begin{table}[htbp]
  \centering
  \caption{三问的递进式建模路线}
  \label{tab:roadmap}
  \small
  \begin{tabularx}{\textwidth}{@{}l X X X@{}}
    \toprule
    问题 & 基线 & 改进 & 最终 \\
    \midrule
    问题一 & 逐维均值/岭回归 & 缺失指示 + 风格先验的多输出回归 & 五模型等权融合 + 保序校准 \\
    问题二 & 边际相关排序 & 图拉普拉斯平滑的重要性 & 多方法集成 + bootstrap 稳定性选择 \\
    问题三 & 风格均值/协方差高斯采样 & 条件高斯 copula & OT + style cap 缺陷投影 \\
    \bottomrule
  \end{tabularx}
\end{table}

\figauto[1.0\textwidth]{fig_roadmap.pdf}{本文总体建模路线}{fig:roadmap}

全文安排如下：第~\ref{sec:preprocess} 节给出数据对齐、缺失处理、变换核验与符号约定；第~\ref{sec:q1} 节至第~\ref{sec:q3} 节分别回答问题一至问题三；第~\ref{sec:conclusion} 节总结结论、建议与局限；附录给出代码结构与复现步骤。
